% GENERATED FILE. Edit canonical lessons, then run npm run generate:books.
% canonical-source-hash: fnv1a64:33cebae41bf24ae6
% canonical-lessons: JA-C355-tabun, JA-C355-kitto, JA-C355-osoraku, JA-C355-kanousei, JA-C355-hazureru

\chapter{How Likely Is It?}
\label{ch:ja-how-likely-is-it}
\hlchaptermodality{355}

\begin{chapteropening}
\textbf{By the end of this chapter:} \emph{I can say how likely something is to happen.}

Complete the last lesson of chapter 355: I can say how likely something is to happen.
\end{chapteropening}

\section[\texorpdfstring{\ja{たぶん} (tabun)}{たぶん (tabun)}]{\ja{たぶん} (tabun) — perhaps; probably}

\label{lesson:JA-C355-tabun}



\begin{quote}
Before the new one: say the Japanese for a parasol, then the Japanese for in advance.
\end{quote}

\subsection*{You'll want to know: \ja{たぶん}}
\textbf{\ja{たぶん}} — \emph{tabun} — \textquotedblleft{}perhaps; probably\textquotedblright{}.

\begin{quote}
\textbf{\ja{あめが} \ja{ふったら、たぶん} \ja{ちちは} \ja{うちに} \ja{いるでしょう。}} — \emph{ame ga futtara, tabun chichi wa uchi ni iru deshō} — \textquotedblleft{}If it rains, my father will probably stay at home.\textquotedblright{}
\end{quote}

\begin{cousinweb}
\textbf{\ja{たぶん}} (\emph{tabun}) is Sino-Japanese (kango): two Chinese characters read with their on'yomi, readings borrowed from Chinese. They are \emph{ta}, \textquotedblleft{}many\textquotedblright{} (Middle Chinese \emph{ta}) and \emph{bun}, \textquotedblleft{}part\textquotedblright{} (\emph{pjun}) (KANJIDIC2; Middle Chinese in Baxter's transcription).
\end{cousinweb}

\begin{grammarlens}[title={What you've built}]
One more word for saying what happens if.
\end{grammarlens}

\subsection*{Your turn}
\begin{itemize}
\raggedright
  \item \emph{Say it:} \emph{tabun}
  \item \emph{Say it:} \emph{tabun}, once more
  \item \emph{Say it:} \emph{ame ga futtara, tabun chichi wa uchi ni iru deshō} — If it rains, my father will probably stay at home.
  \item \emph{Recall:} say the Japanese for a parasol, then the Japanese for in advance, then say \emph{tabun} again
\end{itemize}

\begin{tcolorbox}[breakable,colback=teal!4,colframe=teal!35!black,title={Before you move on}]
What does \ja{たぶん} mean? (\textquotedblleft{}Perhaps; probably\textquotedblright{}.) Say it once more.
\end{tcolorbox}
\section[\texorpdfstring{\ja{きっと} (kitto)}{きっと (kitto)}]{\ja{きっと} (kitto) — surely; certainly}

\label{lesson:JA-C355-kitto}



\begin{quote}
Before the new one: say the Japanese for in advance, then the Japanese for perhaps.
\end{quote}

\subsection*{You'll want to know: \ja{きっと}}
\textbf{\ja{きっと}} — \emph{kitto} — \textquotedblleft{}surely; certainly\textquotedblright{}.

\begin{quote}
\textbf{\ja{まいにち} \ja{れんしゅうしたら、きっと} \ja{かつでしょう。}} — \emph{mainichi renshū shitara, kitto katsu deshō} — \textquotedblleft{}If they practise every day, they will surely win.\textquotedblright{}
\end{quote}

\begin{cousinweb}
\textbf{\ja{きっと}} (\emph{kitto}) is a native mimetic word, one of the Japanese words that picture a manner or a feeling in sound; JMdict marks it as mimetic. It is inherited, not borrowed, and its older origin is not settled (JMdict).
\end{cousinweb}

\begin{grammarlens}[title={What you've built}]
One more word for saying what happens if.
\end{grammarlens}

\subsection*{Your turn}
\begin{itemize}
\raggedright
  \item \emph{Say it:} \emph{kitto}
  \item \emph{Say it:} \emph{kitto}, once more
  \item \emph{Say it:} \emph{mainichi renshū shitara, kitto katsu deshō} — If they practise every day, they will surely win.
  \item \emph{Recall:} say the Japanese for in advance, then the Japanese for perhaps, then say \emph{kitto} again
\end{itemize}

\begin{tcolorbox}[breakable,colback=teal!4,colframe=teal!35!black,title={Before you move on}]
What does \ja{きっと} mean? (\textquotedblleft{}Surely; certainly\textquotedblright{}.) Say it once more.
\end{tcolorbox}
\section[\texorpdfstring{\ja{おそらく} (osoraku)}{おそらく (osoraku)}]{\ja{おそらく} (osoraku) — very likely; in all likelihood}

\label{lesson:JA-C355-osoraku}



\begin{quote}
Before the new one: say the Japanese for perhaps, then the Japanese for surely.
\end{quote}

\subsection*{You'll want to know: \ja{おそらく}}
\textbf{\ja{おそらく}} — \emph{osoraku} — \textquotedblleft{}very likely; in all likelihood\textquotedblright{}.

\begin{quote}
\textbf{\ja{たいふうが} \ja{きたら、おそらく} \ja{でんしゃは} \ja{とまるでしょう。}} — \emph{taifū ga kitara, osoraku densha wa tomaru deshō} — \textquotedblleft{}If the typhoon comes, the trains will very likely stop.\textquotedblright{}
\end{quote}

\begin{cousinweb}
\textbf{\ja{おそらく}} (\emph{osoraku}) is a native Japanese word (wago), inherited rather than borrowed from Chinese, and its older origin is not settled. JMdict notes that it is usually written in kana (JMdict).
\end{cousinweb}

\begin{grammarlens}[title={What you've built}]
One more word for saying what happens if.
\end{grammarlens}

\subsection*{Your turn}
\begin{itemize}
\raggedright
  \item \emph{Say it:} \emph{osoraku}
  \item \emph{Say it:} \emph{osoraku}, once more
  \item \emph{Say it:} \emph{taifū ga kitara, osoraku densha wa tomaru deshō} — If the typhoon comes, the trains will very likely stop.
  \item \emph{Recall:} say the Japanese for perhaps, then the Japanese for surely, then say \emph{osoraku} again
\end{itemize}

\begin{tcolorbox}[breakable,colback=teal!4,colframe=teal!35!black,title={Before you move on}]
What does \ja{おそらく} mean? (\textquotedblleft{}Very likely; in all likelihood\textquotedblright{}.) Say it once more.
\end{tcolorbox}
\section[\texorpdfstring{\ja{かのうせい} (kanōsei)}{かのうせい (kanōsei)}]{\ja{かのうせい} (kanōsei) — a possibility; a chance}

\label{lesson:JA-C355-kanousei}



\begin{quote}
Before the new one: say the Japanese for surely, then the Japanese for very likely.
\end{quote}

\subsection*{You'll want to know: \ja{かのうせい}}
\textbf{\ja{かのうせい}} — \emph{kanōsei} — \textquotedblleft{}a possibility; a chance\textquotedblright{}.

\begin{quote}
\textbf{\ja{くもったら、あめの} \ja{かのうせいが} \ja{あるでしょう。}} — \emph{kumottara, ame no kanōsei ga aru deshō} — \textquotedblleft{}If it clouds over, there will be a chance of rain.\textquotedblright{}
\end{quote}

\begin{cousinweb}
\textbf{\ja{かのうせい}} (\emph{kanōsei}) is Sino-Japanese (kango): three Chinese characters read with their on'yomi, readings borrowed from Chinese. They are \emph{ka}, \textquotedblleft{}can\textquotedblright{} (Middle Chinese \emph{khaX}), \emph{nō}, \textquotedblleft{}ability\textquotedblright{} (\emph{nong}) and \emph{sei}, \textquotedblleft{}nature\textquotedblright{} (\emph{sjengH}) (KANJIDIC2; Middle Chinese in Baxter's transcription).
\end{cousinweb}

\begin{grammarlens}[title={What you've built}]
One more word for saying what happens if.
\end{grammarlens}

\subsection*{Your turn}
\begin{itemize}
\raggedright
  \item \emph{Say it:} \emph{kanōsei}
  \item \emph{Say it:} \emph{kanōsei}, once more
  \item \emph{Say it:} \emph{kumottara, ame no kanōsei ga aru deshō} — If it clouds over, there will be a chance of rain.
  \item \emph{Recall:} say the Japanese for surely, then the Japanese for very likely, then say \emph{kanōsei} again
\end{itemize}

\begin{tcolorbox}[breakable,colback=teal!4,colframe=teal!35!black,title={Before you move on}]
What does \ja{かのうせい} mean? (\textquotedblleft{}A possibility; a chance\textquotedblright{}.) Say it once more.
\end{tcolorbox}
\section[\texorpdfstring{\ja{はずれる} (hazureru)}{はずれる (hazureru)}]{\ja{はずれる} (hazureru) — to miss the mark; to turn out wrong; to come off}

\label{lesson:JA-C355-hazureru}



\begin{quote}
Before the new one: say the Japanese for very likely, then the Japanese for a possibility.
\end{quote}

\subsection*{You'll want to know: \ja{はずれる}}
\textbf{\ja{はずれる}} — \emph{hazureru} — \textquotedblleft{}to miss the mark; to turn out wrong; to come off\textquotedblright{}.

\begin{quote}
\textbf{\ja{よそうが} \ja{はずれたら、みんな} \ja{こまるでしょう。}} — \emph{yosō ga hazuretara, minna komaru deshō} — \textquotedblleft{}If the forecast turns out wrong, everyone will be in trouble.\textquotedblright{}
\end{quote}

Polite \textbf{\ja{はずれます}}, \emph{hazuremasu}; an ichidan verb. Past (polite): \textbf{\ja{よそうが} \ja{はずれました。}}, \emph{yosō ga hazuremashita}.

\begin{cousinweb}
\textbf{\ja{はずれる}} (\emph{hazureru}) is a native Japanese word (wago), inherited rather than borrowed from Chinese: the intransitive partner of \emph{hazusu}, \textquotedblleft{}to take off\textquotedblright{}, and its kanji was attached later as a kun'yomi (JMdict, KANJIDIC2).
\end{cousinweb}

\begin{grammarlens}[title={What you've built}]
One more word for saying what happens if.
\end{grammarlens}

\subsection*{Your turn}
\begin{itemize}
\raggedright
  \item \emph{Say it:} \emph{hazureru}
  \item \emph{Say it:} \emph{hazureru}, once more
  \item \emph{Say it:} \emph{yosō ga hazuretara, minna komaru deshō} — If the forecast turns out wrong, everyone will be in trouble.
  \item \emph{Recall:} say the Japanese for very likely, then the Japanese for a possibility, then say \emph{hazureru} again
\end{itemize}

\begin{tcolorbox}[breakable,colback=teal!4,colframe=teal!35!black,title={Before you move on}]
What does \ja{はずれる} mean? (\textquotedblleft{}To miss the mark; to turn out wrong; to come off\textquotedblright{}.) Say it once more.
\end{tcolorbox}
